\begin{savequote}
\begin{verbatim}
 _______________________________
( Don't look back, the lemmings )
( are gaining.                  )
 -------------------------------
       o   ,__,
        o  (oo)____
           (__)    )\
              ||--|| *

\end{verbatim}
\end{savequote}
\chapter
[\CO{[EIN]} Einführung in die Medizin]
{\CC{[EIN]} Einführung in die Medizin}
\label{chp:EIN}

\ChapterIntro
{%
Die Medizin ist die Lehre von Gesundheit und
Krankheit sowie die Kunde der Gesunderhaltung, Heilung und
Linderung.
}
{%
\begin{description}
  \item[Maintainer:] --
  \item[Autoren:] Diverse
  \item[Reviewer:] Standard-Reviewprozess
  \item[Version:] {\DocVersion} ({\DocVersionInfo})
  \item[SHA1:] {\DocGitCommit}
\end{description}

{\TxTeilkapitelAASS}
}



\clearpage % TEMP temp clearpage
\section{Begriffe}

\nocite{BrockhausDrei1Bd2}
\nocite{BrockhausDrei2Bd2}
\nocite{BrockhausDrei3Bd2}
\nocite{BertelsmannVolkslexikon25}
\nocite{Pschy259}

Die Kenntnis von Fachbegriffen in der Medizin und
Sanitätshilfe ist nicht deshalb wichtig, weil die Fachwörter
so gut klingen. Vielmehr dient die Fachsprache dazu, dass
verschiedene Personen miteinander kommunizieren können und
dabei sichergestellt ist, dass sie sich gegenseitig
verstehen.


\subsection{Allgemeine Begriffe}


\Layout{2}{Symbole}{

\begin{multicols}{2}
\begin{description}
  \item[\Female] Frau, weiblich
  \item[\Male] Mann, männlich
  \item[\Min] Minute 
  \item[\Sek] Sekunde
\end{description}
\end{multicols}
}{}{}{}{}


\Layout{2}{Klinik -- Diagnose -- Therapie}
{

\begin{multicols}{2}
\begin{description}
  \item[Symptome]
\index{Symptom}
Krankheitserscheinungen

\item[Syndrom]
\index{Syndrom}
definierte Kombination von Symptomen

\item[Komplikationen]
\index{Komplikation}
möglicherweise auftretende Schwierigkeiten
\item[Therapie] \index{Therapie} Ma{\ss}nahmen zur
Heilung oder Linderung von Krankheiten und Symptomen

\item[Indikation] Grund zur Anwendung eines bestimmten Verfahrens. 
\item[Kontraindikation] Grund zur \E{Nicht}anwendung eines bestimmten
Verfahrens. Es gibt \E{relative} und \E{absolute} Kontraindikationen. Relative
Kontraindikationen treffen nur unter bestimmten Umständen oder in gewissem
Umfang zu, absolute Kontraindikationen verbieten das Verfahren praktisch immer. 
\item[Diagnose] Zuordnung von Beschwerden und Symptomen zu
einem definierten Krankheitsbild. Es gibt verschiedene Arten
und Funktionen von Diagnosen, siehe
\FullPrettyRef{topic:diagnose}.
\item[Diagnostik]\index{Diagnostik}

Als Diagnostik versteht man alle Bemühungen um eine Diagnose zu stellen, egal
ob man sich der eignen Sinne oder komplizierter Geräte
bedient.

\item[Befund]\index{Befund}

Jede diagnostische Maßnahme hat (hoffentlich) ein Ergebnis, welches man als
\T{Befund} bezeichnet. 
\end{description}
\end{multicols}
}
{}
{}
{}
{}




\Layout{2}{Vorsilben}{

\begin{multicols}{2}
\begin{description}
  \item[Hyper-] hoch, oben
  \item[Hyperglykämie] hoher Blutzucker
  \item[Hypertonie] Bluthochdruck
  \item[Hypo-] tief, unten
  \item[Hypoglykämie] niedriger Blutzucker
  \item[Hypotonie] niedriger Blutdruck
  \item[Brady-] langsam
  \item[Bradykardie] langsamer Herzschlag
  \item[Bradypnoe] langsame Atmung
  \item[Tachy-] schnell
  \item[Tachykardie] schneller Herzschlag
  \item[Tachypnoe] schnelle Atmung
\end{description}
\end{multicols}
}{}{}{}{}

\Layout{2}{Beschreibung von Krankheiten und Symptomen}
{\begin{multicols}{2}
\begin{description}
    \item[akut] plötzlich; rascher Handlungsbedarf
    \item[chronisch] dauerhaft, wiederkehrend
    \item[\Z{rezidivierend}] \Z{wiederkommend, neuerlich}

\end{description}
\end{multicols}}
{}
{}
{}
{}


\Layout{2}{Stoffe}{
\begin{multicols}{2}
\begin{description}
  \item[Aminosäuren] Bausteine von Proteinen.
  \item[Base] \Lang{Syn.} \I{Lauge}. Stoff, welcher
  Protonen aufnehmen kann.\ZFootnote{Definition nach Brønstedt.}
  \Z{Sie reagieren mit Säuren und bilden Salze.}
  \cite{SchmuckChemie1}
  \item[Enzym] \label{topic:enzym}Hilfsstoff, welcher eine chemische Reaktion
  (Umwandlung von Stoffen in andere Stoffe) begünstigt.
  \item[Protein] Eiweiß. Proteine werden aus Aminosäuren
  zusammengesetzt.
  \item[Säure] Stoff, welcher Protonen abgeben kann. \Z{Sie
  reagieren mit Basen und bilden Salze.}
  \cite{SchmuckChemie1}.
  
\end{description}
\end{multicols}
}{}{}{}{}


\Layout{2}{Medizinische Fachrichtungen und Fachabteilungen}{
\begin{multicols}{2}
\begin{description}
  \item[Anatomie] Lehre vom Bau der
  Körperteile~\cite{Pschy259}.
  \item[Chirurgie] 
  \item[Gynäkologie] Frauenheilkunde
  \item[Hygiene]
  \item[\Z{ICU}] \Z{\Lang{engl.} \emph{Intensive Care
  Unit}; Intensivstation}
  \item[\Z{IBS}] \Z{Intensivbehandlungsstation,
  Intensivstation}
  \item[Innere Medizin]
  \item[Kardiologie] Herzheilkunde, Spezialgebiet der inneren
  Medizin
  \item[Neurologie] Fachgebiet der Medizin, welches sich 
  mit der Erforschung, Diagnostik und Behandlung von
  Erkrankungen des \E{Nervensystems} und der Muskulatur
  befasst.
  \item[Nuklearmedizin] 
  \item[Onkologie] 
  \item[Pädiatrie] Kinderheilkunde, Patienten bis zur
  Vollendung des 18. Lebensjahres
  \item[PCI] \Lang{Abkz.} \emph{\Z{Perkutane
  Coronar-Intervention}}; Herzkatheterbehandlung, siehe PTCA
  \item[Physikalische Medizin]
  \item[Psychiatrie]
  \item[PTCA] \Lang{Abkz.} \emph{\Z{Perkutane Transluminale
  Coronal-Angiographie}}; Herzkatheteruntersuchung
  \item[Pulmologie] Lungenheilkunde, Spezialgebiet der
  inneren Medizin
  \item[Radiologie]
  \item[Stroke Unit] Spezialstation zur Behandlung von
  \E{Schlaganfällen}
  \item[Toxikologie] Lehre der Vergiftungen und Giftstoffe
  \item[Traumatologie] Lehre und Fachgebiet der Medizin,
  welches sich mit Unfällen und Verletzungen beschäftigt
  \item[Urologie] Heilkunde des harnableitenden Systems
\end{description}
\end{multicols}
}{}{}{}{}


\Layout{2}{Sonstige Begriffe}{
\begin{multicols}{2}
\begin{description}
  \item[Bolus] \begin{inparaenum}
    \item \enquote{Happen}; 
    \item Stoßweise
  Verabreichung eines Medikaments (intravenös)
  \end{inparaenum}
  \item[Cave] \enquote{Achtung!}, \enquote{Hüte Dich!}
  \item[Obstruktion] Verlegung, Verengung. z.\,B. im Rahmen
  der \Abk{COPD }(\prettyref{topic:copd}).
  \opt{STATUS-EXPERIMENTAL}{\item[Organ] \TODO{GABS}{yyyy-mm-dd}{EIN /
  Begriffe: \enquote{Organ} muss noch erklärt werden.}{}}
  \item[Pathologie, pathologisch, Patho-] Lehre von den
  Krankheiten. Als pathologisch wird alles Krankhafte
  bezeichnet. Die Vorsilbe \enquote{Patho} weist auf etwas
  Abnormales, Krankhaftes hin.
  \item[Zyanose] Blaufärbung der Haut aufgrund eines mangelnden
  Sauerstoffgehaltes des Blutes. Betroffen sind insbesonders die Finger, Lippen
  und das Gesicht.
\end{description}
\end{multicols}
}{}{}{}{}






\subsection{Diagnose} 
\index{Diagnose}
\label{topic:diagnose}
Der Begriff Diagnose wird dauernd verwendet, dennoch ist er
nicht einfach zu erklären. Eine Diagnose gibt einem Zustand
eine Bezeichnung, das heißt einem Beschwerdebild wird ein
definiertes Krankheitsbild zugeordnet.

\begin{center}\FormatTable
\noindent\begin{tabularx}{\linewidth}{XcX}\toprule
Beschwerden und Symptome des Patienten
&
{\color{darkBlue}\sffamily\bfseries $\rightarrow$~DIAGNOSE~$\leftarrow$}
&
Definiertes Krankheitsbild
\\\addlinespace
Patient
&
{\color{darkBlue}\sffamily\bfseries $\rightarrow$~DIAGNOSE~$\leftarrow$}
&
Lehrbuch
\\\addlinespace
Realität
&
{\color{darkBlue}\sffamily\bfseries $\rightarrow$~DIAGNOSE~$\leftarrow$}
&
Lehre
\\\bottomrule
\end{tabularx}
\end{center}


\Fazit{Die Diagnose ordnet die Beschwerden des Patienten einem definierten
Krankheitsbild zu.}

\Layout{0}{Arten von Diagnosen}
{Diese Zuordnung kann einfach sein, sehr oft ist die Diagnosestellung
aber unsicher. Um die Unsicherheit dieser Diagnosen
anzuzeigen haben sich verschiedene Begriffe eingebürgert: 

\begin{itemize}
    \item \T{Verdachtsdiagnose}: 
    Die Diagnose ist unsicher, man hat zwar den Verdacht in Richtung  eines
    bestimmten Krankheitsbildes, andere ähnliche Krankheitsbilder können  aber
    nicht mit hinreichender Sicherheit ausgeschlossen werden. Die endgültige 
    Klärung überlässt man jemand anderen\footnote{z.\,B. dem Krankenhaus}. \par 
    Sanitäter erstellen i.\,d.\,R. nur Verdachtsdiagnosen.
     
    \item \T{Arbeitsdiagnose}: 
    Arbeitsdiagnosen sind Verdachtsdiagnosen, bei denen eine endgültige Klärung
    nicht  möglich ist. Für die Behandlung wird die wahrscheinlichste
    oder gefährlichste Diagnose ausgewählt um \emph{``am Patienten arbeiten zu
    können''.}
    \item \T{Notfalldiagnose}: Eine
    Notfalldiagnose ist im Prinzip eine Arbeitsdiagnose. Der Patient ist dabei 
    in einem Zustand bei dem \EE{sofortiges Handeln} notwendig ist und die 
    eigentliche Diagnose, die zu dieser Situation geführt hat, in dem Moment 
    unwichtig ist. \par Ein Beispiel hierfür sind die Bewusstlosigkeit und der
    Kreislaufstillstand. Die Ursachen dafür sind erstmal unwichtig, mit der
    entsprechenden Behandlung der Notfalldiagnose muss sofort begonnen werden.
    \item \T{Status post},  \T{Zustand nach}:
    \index{St.\,p.|see{Status post}}
    \index{Z.\,n.|see{Zustand nach}} \Lang{Abk.}
    \E{'St.\,p.'}, \E{'Z.\,n.'}. Bezeichnet eine \enquote{alte
    Diagnose}, d.\,h. eine Erkrankung oder eine Behandlung
    die einmal durchgemacht wurde. Z.\,B.:
    \begin{itemize}
        \item \Tt{Z.\,n. Blinddarm-Operation}: Der Patient
        hatte einmal eine Blinddarm-Operation.
        \item \Tt{Z.\,n. Herzinfarkt 2/2006}: Der Patient
        hatte im Februar 2006 einen Herzinfarkt.
    \end{itemize}
    \emph{``Umgangssprachlich''} wird oft auch der Umstand der zu einer
    Diagnose geführt hat mit \emph{``Z.\,n.''} angegeben. 
    \begin{itemize}
        \item \Tt{V.\,a. Schenkelhalsfraktur re., Z.\,n. Sturz}: Der
        Patient ist gestürzt und hat deshalb möglicherweise eine
        Schenkelhalsfraktur rechts.
    \end{itemize}
\end{itemize}}
{
\begin{itemize}
    \item {Verdachtsdiagnose}
    \item {Arbeitsdiagnose}
    \item {Notfalldiagnose}
    \item {Status post}, {Zustand nach}
\end{itemize}}
{}
{}
{}




\Layout{0}{Diagnosecodes}
{Zu statistischen und organisatorischen Zwecken werden
Diagnosen häufig zur weiteren Verarbeitung und Auswertung
codiert, d.\,h. die jeweilige Diagnose wird einem
definierten Code zugeordnet. 

Ein weit verbreitetes System
ist ist die \I{Internationale statistische Klassifikation
der Krankheiten und verwandter Gesundheitsprobleme}
(\T{ICD}\footnote{\Lang{engl.} International Statistical
Classification of Diseases and Related Health Problems}),
welche von der Weltgesundheitsorganisation (WHO)
herausgegeben wird. Derzeit aktuell ist die Version 10
(\enquote{\T{ICD-10}}). \Z{Ein ICD-10-Code beginnt mit einem
Buchstaben. Die ersten drei Stellen geben eine grobe
Diagnose an. Stellen 4 und 5 werden von den ersten drei Stellen mit einem Punkt getrennt und erlauben eine genauere
Beschreibung der Diagnose.}

\Z{\EE{Beispiel}:

\begin{itemize}
  \item \EE{\texttt{S50-S59}}: Verletzungen des Ellenbogens
  und des Unterarmes
  \item \EE{\texttt{S52.3-}}: Fraktur des Radiusschaftes
  \item \EE{\texttt{S52.31}}: Fraktur des distalen
  Radiusschaftes mit Luxation des Ulnakopfes
\end{itemize}}}
{\begin{itemize}
  \item ICD-10
\end{itemize}}
{}
{}
{}

\Layout{2}{Differentialdiagnosen (DD)}
{\index{Differentialdiagnose}%
Diagnosen, welche bei den jeweiligen Symptomen auch
möglich, aber nicht so wahrscheinlich sind wie jene
Diagnose für welche man sich \enquote{entscheidet}.

\begin{quote}
Beispiel: Patient mit Kopfschmerzen nach ausgiebigen
Alkoholgenuss am Vortag.
\begin{description}
    \item[Verdachtsdiagnose:] \E{\enquote{Kopfschmerzen}}
    \item[Differentialdiagnosen:]
    Migraine, Hirnblutung, Schlaganfall,
    Hypoglykämie, Intoxikation, \ldots
\end{description}
\end{quote}}
{}
{}
{}
{}






\bigskip

% \clearpage % TEMP temp clearpage
\section{Wichtige Pathomechanismen}


\subsection{Das Gewebe betreffende Störungen}

Durch krankhafte Vorgänge kann es im Gewebe zu
\begin{itemize}
    \item \T[Oedem@Ödem]{Ödemen} (Wasseransammlung im
    Gewebe),
    \item Emphysemen (Luftansammlung im Gewebe),
    \item Entzündungen (\prettyref{topic:entzuendung}),
    \item Tumoren (\prettyref{topic:tumor}) oder
    \item Nekrosen (Absterben von Zellen bzw. Gewebe)
\end{itemize}
kommen. Die lezten drei Punkte werden im Folgenden näher
erklärt.

\subsection{Entzündung}
\index{Entzündung}
\label{topic:entzuendung}
\HeadingSub{\Z{\protect\Lang{engl., lat.} inflammation.}}

% \DefinitionExp{Entzuendung}
\Definition{Reaktion des Körpers welche charakterisiert
ist durch \begin{inparaenum}
    \item \EE{Rötung},
    \item \EE{Schwellung},
    \item \EE{Überwärmung},
    \item \EE{Schmerz} und
    \item \EE{Funktionsverlust}.
\end{inparaenum}}

% \Fazit{Entzündung: \EE{Rötung}, \EE{Schwellung},
% \EE{Überwärmung}, \EE{Schmerz}, \EE{Funktionsverlust}.}

\paragraph{Auslöser}

\begin{itemize}
    \item Fremdkörper
    \item Krankheitserreger
    \item Mechanische Beanspruchung
    \item Chemische Substanzen
    \item \ldots{}
\end{itemize}

\paragraph{Anhängsel
{\textquotedblleft{}\T{-itis}{\textquotedblright}}}\index{itis@-itis}
Gastritis ${\rightarrow}$
{\quotedblbase}Magenentzündung{\textquotedblleft}, Hepatitis
${\rightarrow}$ Leberentzündung, Pankreatitis ${\rightarrow}$
Entzündung der Bauchspeicheldrüse, Meningitis ${\rightarrow}$
Hirnhautentzündung, Rhinitis ${\rightarrow}$
{\quotedblbase}Schnupfen{\textquotedblleft}, Otitis ${\rightarrow}$
Ohrenentzündung, ...


\subsection{Ischämie}
\label{topic:ischaemie}

\T{Ischämie}  bezeichnet die akute Unterversorgung eines
Gewebes mit sauerstoffreichem Blut infolge der Verstopfung des dafür
zuständigen Blutgefäßes. Eine Ischämie kann grundsätzlich
überall im Körper vorkommen. Es kommt zu Funktionsausfällen
des Gewebes bishin zum Gewebstod. Schermzen sind ein
typisches Symptom. Häufig betroffene Gewebe sind:

\begin{itemize}
  \item Herzmuskel (Herzinfarkt,
  \FullPrettyRef{topic:herzinfarkt})
  \item Hirn (trockener Insult,
  \FullPrettyRef{topic:insult})
  \item Darm (Mesenterialinfarkt, \FullPrettyRef{topic:mesenterialinfarkt})
  \item Extremitäten (Periphere arterielle
  Verschlusskrankheit, arterieller peripherer
  Gefäßverschluss, \FullPrettyRef{topic:arterieller-gefaessverschluss})
\end{itemize}


\subsection{Tumore -- Neubildungen -- Krebs}
\label{topic:tumor}


Die meisten Zellen des menschlichen Körper teilen sich
regelmäßig. Diese Form der Zellvermehrung ist für das
Überleben des Menschen notwendig. Eine Kontrolle dieses
Vorgangs ist jedoch erforderlich um zu gewährleisten dass
die Körperzellen nicht wild zu wuchern beginnen. Der Körper
verfügt über verschiedene Mechanismen um die Zellteilung im
richtigen Maß zu halten. 

Wenn diese Mechanismen versagen, kommt es zu einer
\E{Entartung} und \E{Wucherung}, meist in Form einer
\EE{Raumforderung} und Schwellung (\T{Tumor}\footnote{Der
Begriff \E{\enquote{Tumor}} bedeutet ganz allgemein \enquote{Schwellung}. Umgangssprachlich wird er oft
(fälschlicherweise) mit \enquote{Krebs} gleichgesetzt.}). Man
spricht dann von einer Neubildung \Z{(Neoplasie)}, in
Arztbriefen oft mit \enquote{{\ttfamily N.}} abgekürzt. Der
medizinische Fachbereich der sich mit
Neubildungen (Krebserkrankungen) befasst ist die
\E{Onkologie}.

Man unterscheidet zwischen gutartigen \Z{(benignen)} und
bösartigen \Z{(malignen)} Neubildungen. Gutartige
Neubildungen \enquote{bleiben wo sie sind} und wachsen vor sich
hin. Bösartige Neubildungen bilden Absiedelungen,
sog. \E{Metastasen}, welche z.\,B. über die Blutbahn an
andere Stellen des Körper gelangen und dort zu wuchern
beginnen. 

Grundsätzlich können alle teilungsfähigen Zellen entarten.
Dies kann spontan geschehen oder durch andere Faktoren
(radioaktive Strahlung, chemische Substanzen) begünstigt
bzw. ausgelöst werden.

\Z{Je nach Zellart unterscheidet man bei bösartigen
Wucherungen zwischen Karzinomen, Sarkomen und Blastomen.}

\begin{table}[H]\FormatTable
\Captionof{table}{Beispiele: Häufige Neubildungen.}

\begin{tabularx}{\linewidth}{XXl}
\toprule\addlinespace
\EE{Schlagwort}
&
\EE{Beispiel}
&
\TabHeading{\Z{Im Arztbrief}}
\\\addlinespace\midrule\addlinespace
Lungenkrebs
&
Bronchialkarzinom
&
{\ttfamily\Z{N. bronchii}}
\\\addlinespace
Brustkrebs
&
Mammakarzinom
&
{\ttfamily\Z{N. mammae}}
\\\addlinespace
Darmkrebs
&
Rektumkarzinom
&
{\ttfamily\Z{N. recti}}
\\\addlinespace
Blutkrebs
&
Leukämie, versch. Arten
&
{\ttfamily\Z{ALL, AML, CML, \ldots}}
\\\addlinespace\bottomrule
\end{tabularx}
\end{table}

\paragraph{Was ist das Problem?}

Zusammengefasst ergeben sich drei wesentliche Probleme:

\begin{enumerate}
    \item \EE{Neubildungen brauchen Raum.} Dieser Raum ist
    nicht endlos vorhanden. Je nach Lage können wichtige
    Strukturen beengt oder behindert werden (Hirn,
    Luftröhre, \ldots).
    \item \EE{Funktionsverlust des gesunden Gewebes.}
    Durch den Befall mit (nutzlosen) Metastasen wird das
    gesunde Gewebe soweit durchwachsen und geschwächt, dass
    es seine Funktion nicht mehr richtig ausführen kann
    (z.\,B. Knochenbefall).
    \item \EE{Schmerz.} Der Befall verursacht oft starke
    Schmerzen. Krebspatienten benötigen oft eine
    dauerhafte, starke Schmerztherapie.
\end{enumerate}

\paragraph{Therapie}

Der Therapieerfolg ist sehr stark von der genauen Art der
Neubildung und vom jeweiligen Patienten abhängig und reicht
von sehr gut bis nicht vorhanden. 

Im wesentlichen kommen bei der Krebstherapie folgende
Methoden zum Einsatz:

\begin{itemize}
    \item \E{Chemotherapie}: Einsatz von Medikamenten welche
    die Zellteilung unterdrücken bzw. Zellen absterben
    lassen.
    \item \E{Bestrahlung}: zur Abtötung oder Vorbeugung von
    Geschwülsten.
    \item \E{Operation}: Chirurgische Entfernung der
    Neubildung(en).
\end{itemize}

Diese Therapien sind für den Körper sehr anstrengend und
haben viele Nebenwirkungen, wie z.\,B. Haarausfall, Übelkeit,
Erbrechen, \EE{Schwächung des Immunsystems}, Nervenschäden
u.\,v.\,a.\,m.

% \subsection{Ischämie}

% TODO INHALT-NEU: Ischämie 

% \section{Reserviert}


\cite{UlfigHistologie1,GantenTumorerkr98,Uicc3}

\bigskip

% \clearpage % TEMP temp clearpage
\section{Dokumentation -- Von der medizinischen Seite
betrachtet}
\index{Dokumentation!medizinische Aspekte}
\label{topic:dokumentation-medizinisch}

\emph{Vgl. auch \prettyref{topic:dokumentation-rechtlich}}

\Layout{60}{Dokumentation -- Oder: Damit sich auch ein
Anderer auskennt \ldots} {Die medizinische Dokumentation muss eine Geschichte erzählen. Eine
Geschichte die ein Dritter, unbeteiligter, verstehen
\E{muss}. Dieser unbeteiligte Dritte soll allein anhand der
Schilderung in Ihrer Dokumentation fähig sein eine
Verdachtsdiagnose zu stellen. Weiters sollte er
nachvollziehen können, was mit dem Patient passiert ist.
Auch die Ereignisse davor oder die Einsatzsituation soll
für denjenigen verständlich sein (Unfallhergang,
Wohnsituation bei psychiatrischen Patienten, \ldots)
}{
}
{./images/teaser/affe-sama-ambas-schein.jpg}
{}
{Gabriel}

\bigskip

\Fazit{Die Dokumentation dient sowohl der rechtlichen Absicherung, als auch
der Beschreibung der Situation gegenüber Dritter, welche am Einsatz nicht
beteiligt waren.}

\bigskip
% \clearpage % TEMP temp clearpage
\subsection{Beispiele}

\Layout{2}{Schlechte Dokumentation}
{\par\bigskip\noindent\begin{minipage}{\linewidth}
{\sffamily Diagnose:} {\ttfamily\color{darkBlue} RQW}

{\sffamily Befunde:} {\sffamily RR}
{\FontTypewriter\color{darkBlue} ~~130/70,} {\sffamily HF}
{\FontTypewriter\color{darkBlue} ~~~90} {\sffamily SpO2}
{\FontTypewriter\color{darkBlue} ~~~98\% }

{\sffamily Anmerkungen: }{\FontTypewriter\color{darkBlue} -- --
--}
\end{minipage}}
{}
{}
{}
{}


\Layout{2}{Bessere Dokumentation}
{
\par\bigskip\noindent\begin{minipage}{\linewidth}
{\sffamily Diagnose:} {\FontTypewriter\color{darkBlue} RQW Stirn}

{\sffamily Befunde:} {\sffamily RR}
{\FontTypewriter\color{darkBlue} ~~130/70,} {\sffamily HF}
{\FontTypewriter\color{darkBlue} ~~~90} {\sffamily SpO2}
{\FontTypewriter\color{darkBlue} ~~~98\% }

{\sffamily Anmerkungen:} {\FontTypewriter\color{darkBlue} Pat.
fuhr mit Fahrrad und prallte mit Stirn gegen Ast., kein Sturz. Lt. Begleiter
keine B.losigkeit, Unfallhergang erinnerlich \underbar{Traumacheck:} RQW Stirn,
ca 50\,ml Blutverlust sonst unauffällig. \underbar{Neuro:}
grob neurologisch unauffällig, Bew. klar, Pat zeitl, örtl, zur Situation \& Person orientiert, Pupillen
mittelweit, Reakt. prompt und seitengleich, Kraft\&Gefühl
seitengl. in OE und UE. \underbar{Anamn.:} \underbar{Sympt:} Kopfschmerz seit
Aufprall. \underbar{Krankh.:} keine bek. \underbar{All.}
keine bek. \underbar{Med.:} Aspirin heute früh. letzte
Tetanus-Impf. nicht bek. \underbar{Ereignisse} s.o.
\underbar{Letzte} Mahlzeit: heute Früh 2 Semmeln\A{ KOCH:
? muss sein? GABS: ja, 1) wegen s-kamel, 2) wegen nüchtern } 

\underline{Maßnahmen:} Steriler Verband, nüchtern belassen,
Transport liegend mit erh OK.}
\end{minipage}
}
{}
{}
{}
{}



\Layout{2}{Psychiatrischer und psychosozialer Patient}
{\par\bigskip\noindent\begin{minipage}{\linewidth}
{\sffamily Diagnose:} {\FontTypewriter\color{darkBlue}
Psychiatrische Erkrankung, V.a. Wahnvorstellungen, Unterbringung nach UbG}

{\sffamily Befunde:} {\sffamily RR}
{\FontTypewriter\color{darkBlue} ~~130/70,} {\sffamily HF}
{\FontTypewriter\color{darkBlue} ~~~90} {\sffamily SpO2}
{\FontTypewriter\color{darkBlue} ~~~98\% }

{\sffamily Anmerkungen:} {\FontTypewriter\color{darkBlue} Pat.
sitzend in Wohnung angetr., von Pol. mit Handschellen gefesselt. Pat. warf lt. Polizei Gegstände aus d Fenster
auf Straße, verletzte Passanten ($\rightarrow$ FLO-1). Pat.
gibt an, sich geg. Geheimdienst gewehrt zu haben,
jetzt agitiert und unkooperativ. 

\underbar{Traumacheck:} keine äußeren Anz f Verletzungen,
sonst nicht erhebbar (unkoop.). \underbar{Neuro:} Pat.
ist wach, sonst nicht erhebbar

\underbar{Anamn.:} \underbar{Sympt:} keine Antwort
\underbar{Krankh.:} lt. Spitalbrief Otto-Wagner-Sp. v.
1.4.08 paranoide Schizophrenie \underbar{All.} k.A.
\underbar{Med.:} nicht eruierbar, auf Tisch mehrere
angebr. Packungen Rohypnol \underbar{Ereignisse} s.o.
\underbar{Letzter Spitalsaufenthalt} k.A.

\underbar{Wohnsituation:} heruntergekommen, vermüllt, Pat.
wohnt mit großem Hund, ganze Wohnung verkotet, bestialisch
stinkend. 

\underbar{Maßnahmen:} Einweisung nach UbG, Transport in
Handschellen in Pol-begl. (Viktor 6, DNr. 4711). Protokoll
Amtsarzt ad. Spital.}
\end{minipage}}
{}
{}
{}
{}





% \section{Reserviert}
% \A{Medizinische Fachbereiche}
% 
% \section{Reserviert}
% \A{Medizinische Berufsgruppen}
% 
% \section{Reserviert}

